\chapter[Discusión]{Discusión general}
\label{cp:discusion}
% Extensión estimada: 10--12 páginas, cinco secciones organizadas por argumentos, sin figuras ni tablas nuevas.
% Reestructurado tras la revisión del autor (28/10): las respuestas a PI1 y PI2 pasan al Capítulo 7.
% Fuentes: Capítulos 4 y 5 (en particular sus balances, sec:res_clasif_discusion y sec:res_gen_discusion, y el material
% argumentativo que salió de ellos en la revisión del tutor de los Caps. 4 y 5); Capítulo 2 (estado del arte y
% tab:marco_brechas); Capítulo 3 (restricción de cómputo, diseño de los experimentos); sección "Discusión y
% conclusiones" de Docs/experimentacion/experimentacion.tex (amenazas a la validez y limitaciones), con las cifras
% actualizadas a la partición por lesión.
%
% Feedback del tutor: cada capítulo de resultados cierra con un balance breve, y este capítulo queda para lo
% transversal, lo que ninguno de los dos podría cerrar por sí solo. Las comparaciones con trabajos relacionados se
% integran en las secciones temáticas, y las amenazas a la validez se fusionan con las limitaciones.

Los Capítulos~\ref{cp:res_clasif} y~\ref{cp:res_gen} presentaron los resultados de la clasificación y de la
generación de reportes, y cada uno cerró con un balance de sus hallazgos. Este capítulo los interpreta en conjunto,
organizado por argumentos y no por el orden de los resultados, y propone para cada uno una explicación, marcada como
tal cuando los experimentos no permiten confirmarla. En primer lugar, se discute dónde está lo que un \gls{vlm} aporta a
la clasificación y qué le falta para diagnosticar. A continuación, se analiza por qué los generadores siguen al texto
que reciben más que a la imagen, y qué implica eso para el diseño y la evaluación del \emph{pipeline}. Después, se
interpreta el ruteo por concordancia como una forma de predicción selectiva, con sus alcances y sus límites, y se
discute qué limita el desempeño del clasificador: los datos, las decisiones de evaluación o la escala del cómputo.
Finalmente, el capítulo reúne las condiciones de validez de sus conclusiones y las limitaciones del trabajo. Las
respuestas a las preguntas de investigación, que se apoyan en esta discusión, se enuncian en el
Capítulo~\ref{cp:conclusiones}.

\section{Especialización de la representación y decisión diagnóstica}
\label{sec:disc_representacion}
% Extensión estimada: 2 páginas, cinco párrafos (ficha 6.1).
% Fuentes: sec:res_vlm_controles (tab:res_vlm_controles), sec:res_vlm_dermlip (tab:res_dermlip, tab:res_dermlip_swin),
% tab:res_supervisados, sec:res_gen_sin_anclaje, sec:res_gen_ablaciones, sec:met_generacion; zeng2025mmskin,
% jiang2026multimodal, yan2026dermfmzero. Reemplaza las antiguas 6.1 (PI1) y 6.2.

Los \glspl{vlm} evaluados ocupan tres niveles en la clasificación de \emph{HAM10000}
(Sección~\ref{sec:res_comparacion}). Los generativos sin especialización en dermatología quedan cerca de 1/7 en sensibilidad macro, DermLIP sin ajuste
queda entre ellos y los supervisados, y DermLIP con una cabeza entrenada se acerca a los supervisados. La diferencia
entre esos niveles no es solo de grado, porque cada uno usa el modelo de una manera distinta. Esta sección discute qué
la explica: qué tiene DermLIP que los generativos no tienen, y qué les falta a estos para convertir lo que ven en un
diagnóstico.

Los controles de la Sección~\ref{sec:res_vlm_controles} permiten acotar las explicaciones posibles. Una de las hipótesis examinadas es que el bajo desempeño de los modelos generativos esté relacionado con la modalidad dermatoscópica, posiblemente menos representada durante su preentrenamiento. Para explorar esta posibilidad, se evalúa su comportamiento sobre fotografías clínicas de \emph{SkinCAP}, una modalidad más cercana a las imágenes de uso general. Las fotografías no mejoran su resultado, mientras que DermLIP distingue las clases en las dos modalidades. Además, con la imagen reemplazada por una gris, los generativos asignan todas las imágenes a una misma clase. La escasa variación de la sensibilidad macro al sustituir la imagen neutra por la imagen real sugiere que, bajo las instrucciones y estrategias evaluadas, la información visual tiene una influencia limitada sobre la decisión final del modelo. En conjunto, los resultados sugieren que las representaciones especializadas en dermatología podrían contribuir al mejor desempeño de clasificación. No obstante, las diferencias entre los modelos evaluados impiden establecer que esta sea la única causa de la ventaja observada.

El control de la representación indica dónde está esa especialización. Con las mismas tres cabezas, el
\emph{embedding} de DermLIP supera al de Swin preentrenado en \emph{ImageNet} por entre 0.15 y 0.35 de sensibilidad
macro (Tabla~\ref{tab:res_dermlip_swin}). Más llamativo es que esa representación, congelada y sin haber visto
\emph{HAM10000}, quede cerca de un clasificador ajustado de punta a punta sobre ese conjunto: con una \gls{svm}, su F1
macro es de 0.705, frente a 0.719 de Swin-Tiny, y con un \emph{linear probe} su sensibilidad macro supera a la de
Swin-Tiny, aunque con un F1 macro menor (Tablas~\ref{tab:res_dermlip} y~\ref{tab:res_supervisados}). Esas diferencias
no se compararon estadísticamente, de modo que solo indican que las dos representaciones quedan en un mismo rango. Una hipótesis compatible con estos resultados es que el entrenamiento contrastivo con pares de imágenes y descripciones dermatológicas favorezca representaciones que distingan características relevantes para la clasificación. Sin embargo, comprobar este mecanismo requeriría un análisis específico de la estructura del espacio de \emph{embeddings} y de su relación con los rasgos clínicos.

Mini-VLM, en cambio, extrae información de la imagen, pero, en estos experimentos, no la convierte en una clase. Después del
ajuste con \emph{SkinCaRe}, la descripción de Mini-VLM cambia con cada lesión, y sin la imagen su cobertura de
entidades clínicas cae a la mitad (Sección~\ref{sec:res_gen_ablaciones}). Sin un diagnóstico, sin embargo, el mismo
generador nombra la clase verdadera en menos de un reporte de cada cinco (Sección~\ref{sec:res_gen_sin_anclaje}). Es decir, la descripción del generador cambia con la imagen, mientras que, sin un diagnóstico, el reporte rara vez
nombra la clase verdadera. Una posible explicación es que las dificultades de clasificación no dependan exclusivamente de la representación visual, sino también de cómo el modelo generativo transforma esa información en una decisión entre categorías. Los resultados disponibles no permiten separar de manera concluyente ambas posibilidades.
Otra explicación, compatible con la anterior, es que \emph{SkinCaRe} agrupa las enfermedades en categorías distintas
de las siete de \emph{HAM10000} (Sección~\ref{sec:met_generacion}), de modo que el ajuste nunca le pidió al generador
esa correspondencia. Los experimentos no permiten separar las dos.

En el caso de SkinVL y con la configuración evaluada, la especialización dermatológica no fue suficiente. Ajustado con imágenes y textos
dermatológicos, sin un diagnóstico escribe casi siempre el mismo texto sobre una infección por hongos
(Sección~\ref{sec:res_gen_sin_anclaje}). Sus propios autores le miden a la variante que usa este trabajo, ajustada solo
con los libros de \emph{MM-Skin}, una exactitud del 3.5\,\% al clasificar \emph{HAM10000} sin ejemplos de ese conjunto, frente al
51.3\,\% de la variante que suma conjuntos de diagnóstico reformulados como preguntas~\cite{zeng2025mmskin}. La especialización de un generador parece quedar ligada, entonces, a las tareas con que se ajustó, y la de una representación
contrastiva se traslada a una tarea nueva con solo una cabeza. La literatura revisada es compatible con esta lectura: los generativos abiertos
obtienen exactitudes bajas también en consultas reales~\cite{jiang2026multimodal}, y por la vía contrastiva
DermFM-Zero supera a DermLIP en el protocolo de sus autores~\cite{yan2026dermfmzero}. Por lo tanto, la utilidad de un
\gls{vlm} para clasificar depende de cómo se lo use: como clasificador autónomo, ninguno de los evaluados alcanza a
los supervisados; como codificador de un modelo especializado, se acerca. El \emph{pipeline} aprovecha esa diferencia,
porque usa la representación de DermLIP para recuperar casos parecidos y toma el diagnóstico de un clasificador
supervisado. Queda por explicar por qué el generador, que no decide la clase, adopta sin más la que recibe.

\section{El peso del texto frente a la imagen en los generadores}
\label{sec:disc_texto_imagen}
% Extensión estimada: 2 páginas, cuatro párrafos (ficha 6.2).
% Fuentes: sec:res_vlm_controles, sec:res_gen_sin_anclaje, sec:res_gen_fuentes, sec:res_gen_ablaciones, sec:res_gen_rag,
% sec:res_gen_comparacion, sec:res_gen_pipeline, sec:res_gen_cualitativos, sec:res_gen_estado_arte, sec:met_generacion
% (instrucción de ajuste sin diagnóstico; LoRA de Mini-VLM solo en el codificador visual); rohrbach2018object,
% bai2024hallucination. Reemplaza la antigua 6.3 y la parte de PI2 de la antigua 6.1.

Los Capítulos~\ref{cp:res_clasif} y~\ref{cp:res_gen} evaluaron a los generativos en tareas distintas, pero en las dos
aparece un mismo patrón: en las condiciones evaluadas, cuando el texto de la entrada y la imagen apuntan a respuestas distintas, el generador tiende a seguir al texto. En la clasificación, los controles son compatibles con que el \emph{prompt} fije una respuesta por defecto que la
imagen desplaza muy poco
(Sección~\ref{sec:res_vlm_controles}).

En la generación de reportes, el patrón aparece en tres situaciones. Sin un diagnóstico en la instrucción, SkinVL repite en siete de cada diez reportes un mismo texto ajeno a la imagen (Sección~\ref{sec:res_gen_sin_anclaje}). En la cascada, el diagnóstico afirmado no cambia al quitar la imagen y, con un
diagnóstico equivocado, el reporte lo afirma en lugar de corregirlo (Sección~\ref{sec:res_gen_ablaciones}). Por
último, cuando recibe a la vez el diagnóstico del clasificador y el de los vecinos, el generador sigue al clasificador
aunque los vecinos acierten: en las 72 imágenes en que el clasificador se equivoca y los vecinos aciertan, recupera la clase verdadera en una de cada seis
(Sección~\ref{sec:res_gen_rag}).

En conjunto, estos resultados muestran que, bajo las condiciones evaluadas, los generadores tienden a conservar el diagnóstico proporcionado en la instrucción, incluso cuando este no coincide con la categoría verdadera. Esta tendencia sugiere una fuerte influencia de la información textual sobre la impresión diagnóstica generada, aunque no permite determinar de manera concluyente el peso relativo de cada modalidad. Además, el diagnóstico de la instrucción pesa más que el contexto recuperado, lo que sugiere una jerarquía entre las fuentes del reporte.

El origen de esa influencia no parece estar en el ajuste de este trabajo. La instrucción utilizada durante el ajuste, «Describe this skin lesion.», no proporciona explícitamente un diagnóstico que el modelo deba reproducir (Sección~\ref{sec:met_generacion}). Además, en Mini-VLM el ajuste no llega al modelo de lenguaje, y SkinVL, que no se ajustó en este trabajo,
transcribe el diagnóstico del mismo modo. Una explicación posible es que la tendencia proviene del modelo de lenguaje
de base, que en su ajuste original aprendió a seguir lo que la instrucción afirma, y que la imagen no alcanza a
contradecir. Es una forma del peso de los patrones del lenguaje que la literatura señala como causa de las
alucinaciones~\cite{rohrbach2018object,bai2024hallucination}, aunque aquí no produce un hallazgo inventado, sino
la adopción de un diagnóstico ajeno. La explicación alternativa es que ceder sea lo esperable en un generador que no
distingue las clases por sí mismo (Sección~\ref{sec:disc_representacion}): si el modelo no logra obtener de la imagen información suficiente para distinguir las categorías, puede tender a utilizar el diagnóstico proporcionado en la instrucción como principal referencia para redactar la conclusión. Las dos explicaciones predicen lo mismo en estos experimentos,
y separarlas requeriría un generador que clasifique bien sin anclaje, que ninguno de los evaluados es.

Cualquiera sea su origen, esa influencia condiciona cómo debe evaluarse un sistema que integra un clasificador y un generador. Como el generador transcribe el
diagnóstico, la fidelidad diagnóstica del reporte mejora con el clasificador más que con el generador: cambiar la
fuente de DermLIP sin ajuste a Swin-Tiny agrega 0.132 de consistencia con el mismo generador
(Sección~\ref{sec:res_gen_fuentes}), mientras que ninguno de los tres generadores eleva la consistencia más de unas milésimas por encima de la exactitud de la fuente. Evaluado como un todo, el sistema puede atribuirle al generador una fidelidad que proviene del clasificador, y por eso debe medirse por componente. En los trabajos incluidos en la revisión no se identificó una evaluación que distinguiera, mediante comparaciones equivalentes, la contribución del clasificador y del generador a la consistencia diagnóstica del reporte final (Sección~\ref{sec:res_gen_estado_arte}). La Tabla~\ref{tab:res_gen_aportes} resume esa separación para este trabajo.

En consecuencia, la selección del generador prioriza su capacidad para conservar el diagnóstico recibido y evitar errores adicionales, dado que su contribución principal al sistema es la producción del contenido textual. LLaVA-Med omite el diagnóstico recibido en uno de cada siete
reportes aun con el \emph{oracle}, y en el nivel diferencial deja los candidatos dentro de un razonamiento sin cerrar
(Secciones~\ref{sec:res_gen_fuentes} y~\ref{sec:res_gen_pipeline}). Por eso el \emph{pipeline} usa Mini-VLM, que conserva el diagnóstico como SkinVL, describe tanto
como LLaVA-Med y entra en una GPU de consumo.

En estos experimentos, esa misma influencia impide que el generador funcione como una segunda opinión, y además alcanza a la descripción. Sin la
imagen, cada generador escribe un único texto por clase, una descripción prototípica del diagnóstico que
recibe; con ella, la descripción varía entre lesiones (Sección~\ref{sec:res_gen_ablaciones}). Una lectura posible es
que el reporte parte de ese prototipo y la imagen lo modifica, pero sin contradecirlo, como en el melanoma afirmado
como nevo y descrito con una pigmentación homogénea «with no visible melanocytic infiltration or atypical features» (Sección~\ref{sec:res_gen_pipeline}).
Qué parte de la descripción proviene del prototipo y qué parte de la imagen no se midió. Lo que sí se sigue de los
resultados es que la verosimilitud de un reporte no indica si su diagnóstico es correcto, porque un error del
clasificador llega al reporte acompañado de rasgos que lo justifican. En este diseño, la señal de cuándo desconfiar debe venir, por lo tanto, de fuera del generador, y esa es la función del ruteo, tema de la sección siguiente.

\section{La concordancia como predicción selectiva}
\label{sec:disc_concordancia}
% Extensión estimada: 2 páginas, cinco párrafos (ficha 6.3).
% Fuentes: sec:marco_pred_selectiva, sec:marco_concordancia, sec:met_ruteo; sec:res_gen_concordancia (1 204 afirmados,
% 102 errores afirmados, 230 errores de Swin-Tiny, 248 diferenciales, 51 derivados), sec:res_gen_rag (171 y 72),
% tab:res_dermlip (kNN 0.781), sec:res_gen_pipeline, sec:res_melanoma, sec:res_calibracion, sec:res_gen_estado_arte.
% Cifras derivadas: cobertura 1 204 / 1 503 = 0.80; riesgo 102 / 1 204 = 0.085; error global 230 / 1 503 = 0.153;
% ambos equivocados 230 - 72 = 158; esperable si fueran independientes 0.153 x 0.219 x 1 503 = 50.
% Reemplaza la parte del ruteo de las antiguas 6.3 y 6.4.

El ruteo por concordancia es, en los términos de la Sección~\ref{sec:marco_pred_selectiva}, un clasificador selectivo:
Swin-Tiny propone el diagnóstico, y la concordancia con el \gls{knn} decide si se afirma. Leído así, sus resultados
(Sección~\ref{sec:res_gen_concordancia}) se resumen en dos cifras. El \emph{pipeline} afirma el diagnóstico en el
80\,\% de las imágenes de prueba, que es su cobertura, y entre ellas el clasificador se equivoca en el 8.5\,\%, que es
su riesgo selectivo, frente al 15.3\,\% de error sobre la prueba completa. Es decir, renunciar a afirmar una de cada
cinco imágenes reduce a algo más de la mitad la proporción de diagnósticos afirmados erróneos. Además, a diferencia de
un umbral sobre la probabilidad, la regla no tiene ningún parámetro ajustado sobre la prueba, de modo que esa
reducción no está sobreestimada por la selección.

Que la concordancia funcione depende de que los dos predictores se equivoquen en imágenes distintas. La
Sección~\ref{sec:marco_concordancia} lo planteaba como una hipótesis: un clasificador que ajusta fronteras entre las
clases y un \gls{knn} que vota con los casos más parecidos, sobre una representación aprendida de otra manera, deberían
tener sesgos distintos. Los resultados la respaldan solo en parte. Los predictores difieren en muchas imágenes: el
\gls{knn} acierta 72 en que Swin-Tiny se equivoca, y Swin-Tiny acierta 171 en que el \gls{knn} se equivoca
(Sección~\ref{sec:res_gen_rag}). Esos desacuerdos son los que el ruteo aprovecha, y por eso en el nivel diferencial el
par de candidatos contiene la clase verdadera en el 84.7\,\% de los casos.

Sus errores, sin embargo, coinciden mucho más de lo que coincidirían por azar. Los dos predictores se equivocan a la
vez en 158 imágenes, y en 102 de ellas eligen la misma clase errónea. Si sus errores fueran independientes, con una
tasa de error de 0.153 para Swin-Tiny y de 0.219 para el \gls{knn}, se equivocarían a la vez en unas 50 imágenes, de
modo que la coincidencia observada es unas tres veces mayor. Una explicación posible es que los dos predictores
aprenden de las mismas 6\,447 imágenes de entrenamiento: el clasificador ajusta sus fronteras con ellas y el \gls{knn}
las usa como casos de referencia. Si una lesión de prueba se parece a lesiones de entrenamiento de otra clase, o si
esas lesiones están mal etiquetadas, los dos heredan la misma confusión, por distintas que sean sus representaciones.
Otra explicación, compatible con la anterior, es que una parte de esas imágenes sea ambigua para cualquier modelo,
como lo es la distinción entre melanoma y nevo para los especialistas. En cualquier caso, la heterogeneidad de las
representaciones no basta para que los errores sean independientes. Otras señales de confianza, como la probabilidad de
Swin-Tiny, que está bien calibrada cuando se entrena sin pesos por clase (Sección~\ref{sec:res_calibracion}), no
intervienen en el ruteo, y no se evaluó si detectan una parte de esos errores.

Los errores compartidos se concentran, además, donde el costo clínico es mayor. El \emph{pipeline} se eligió y se
evaluó por la consistencia diagnóstica, que con la cascada coincide con la exactitud del clasificador y está dominada,
por lo tanto, por los nevos, que son dos tercios de la prueba. En el melanoma, la prioridad clínica es la
sensibilidad, y ahí el \emph{pipeline} queda más débil: casi uno de cada cuatro melanomas se afirma como otra clase, la
mayoría como nevo (Sección~\ref{sec:res_gen_concordancia}). El ruteo mitiga parte del problema, porque cuando un
melanoma pasa al nivel diferencial el par de candidatos casi siempre lo incluye, pero no puede mitigar los melanomas
que los dos predictores confunden. Al mismo tiempo, el umbral fijado en validación, que lleva la detección de
Swin-Tiny de tres a cuatro de cada cinco melanomas (Sección~\ref{sec:res_melanoma}), no se usa en el ancla. Es decir,
el \emph{pipeline} optimiza una métrica que trata igual todos los errores, mientras que la clínica no los trata igual,
y la configuración que maximiza la consistencia no tiene por qué ser la preferible en la práctica.

Lo que sí logra el ruteo es convertir parte de los errores en una duda explícita. En el nivel diferencial, el reporte
presenta dos candidatos en lugar de afirmar el diagnóstico equivocado del clasificador, y en la derivación no se emite
reporte: se evitan 20 reportes incorrectos y dos contradictorios, a costa de 29 diagnósticos correctos. El valor de ese
intercambio depende de cuánto cueste en la práctica un error frente a una consulta con el especialista, una relación
que este trabajo no midió. El único trabajo revisado que combina un ensamble con la derivación,
el de \textcite{khan2025ensemble}, no informa cuántos casos deriva ni con qué acierto (Sección~\ref{sec:res_gen_estado_arte}),
de modo que la cobertura y el riesgo selectivo de este trabajo no tienen todavía con qué compararse. En conjunto, el
\emph{pipeline} puede ordenar la atención del especialista, señalando qué casos revisar primero, pero no reemplazarla,
la misma conclusión a la que llegan \textcite{jiang2026multimodal} sobre los \glspl{vlm} que evalúan en consultas
reales.

\section{Datos, evaluación y escala como límites del desempeño}
\label{sec:disc_limites}
% Extensión estimada: 2.5 páginas, cinco párrafos (ficha 6.4).
% Fuentes: sec:res_sup_arquitecturas, sec:res_sup_clases, tab:res_supervisados (ensamble), sec:res_melanoma, sec:res_fuga,
% sec:res_comparacion, sec:res_estado_arte, sec:res_gen_nivel1, sec:res_gen_sin_anclaje, sec:res_gen_estado_arte,
% sec:res_costo_clasificadores, sec:res_costo, sec:res_gen_comparacion, tab:marco_generadores, sec:estado_arte_vlm;
% kim2025skinsavvy2, yao2022single, khan2025ensemble, yan2026dermfmzero. Reemplaza las antiguas 6.5 y 6.6.

Si el reporte hereda el desempeño del clasificador, la pregunta siguiente es qué limita ese desempeño. Entre Swin-Tiny, ConvNeXt-Tiny y ViT-B/16 no se encontraron diferencias estadísticamente respaldadas bajo las condiciones evaluadas (Sección~\ref{sec:res_sup_arquitecturas}). Por lo tanto, los resultados no permiten establecer una ventaja clara de una arquitectura sobre las otras, aunque tampoco demuestran que su desempeño sea equivalente. Además, ViT-B/16 no supera a Swin-Tiny con el triple de parámetros, aunque usa la configuración elegida para este último, y el ensamble de los dos mejores supera la media de cada uno en unas 0.02 de F1 macro, una diferencia del orden del desvío entre ejecuciones (Tabla~\ref{tab:res_supervisados}). Una explicación posible es que la disponibilidad y distribución de las imágenes de entrenamiento limiten las mejoras obtenidas mediante arquitecturas de mayor capacidad. No obstante, los experimentos realizados no permiten determinar el peso relativo de ambos factores.

Además de los factores que pueden influir sobre el desempeño real de los clasificadores, las decisiones adoptadas durante la evaluación condicionan las métricas obtenidas y la posibilidad de compararlas con otros trabajos. Estas decisiones mueven las cifras en magnitudes iguales o mayores que las que separan a las arquitecturas. La partición por imagen sobreestima la sensibilidad macro en unas 0.07 (Sección~\ref{sec:res_fuga}),
una cifra del mismo orden que la distancia entre Swin-Tiny y el trabajo de \textcite{kim2025skinsavvy2}, que divide las
imágenes al azar. En melanoma, la regla de decisión pesa más que el modelo: el umbral fijado en validación lleva la
detección de tres a cuatro de cada cinco melanomas sin cambiar el \gls{auc} (Sección~\ref{sec:res_melanoma}), mientras
que los trabajos revisados operan en puntos que van desde la clase más probable hasta una sensibilidad cercana a la
total. Y el rango de la sensibilidad entre ejecuciones supera una décima en las clases minoritarias, de modo que un resultado de una
sola ejecución puede quedar arriba o abajo de otro por la variación entre entrenamientos. Por eso una comparación con la literatura es interpretable cuando se conocen la partición, la regla de decisión y la cantidad de ejecuciones, y por eso este trabajo fija el
umbral en validación, a diferencia de un modelo elegido sobre el mismo conjunto de prueba en que se informa, como el de
\textcite{yao2022single} (Sección~\ref{sec:res_estado_arte}).

La elección de la métrica también puede llevar a conclusiones distintas, especialmente en conjuntos desbalanceados. McNemar compara las imágenes que acierta cada modelo, y el \emph{bootstrap} pareado compara su sensibilidad macro. Según McNemar, DenseNet121 acierta más imágenes que DermLIP, 250 frente a 85 discordantes. Según el \emph{bootstrap}, en cambio, su sensibilidad macro no se distingue de la de DermLIP (Sección~\ref{sec:res_comparacion}). Con clases desbalanceadas, el total de aciertos queda dominado por las clases mayoritarias, y por eso la comparación tiene que aplicarse a una métrica macro.

Una dificultad similar aparece en la evaluación de los reportes generados, aunque las métricas empleadas miden propiedades diferentes. Las medidas de similitud textual permiten comparar los reportes con referencias, mientras que la consistencia diagnóstica identifica la presencia de la categoría verdadera. Ninguna de estas medidas, por sí sola, permite establecer la fidelidad clínica completa del texto. Las métricas basadas en referencias son sensibles a las diferencias léxicas y semánticas entre el reporte generado y el texto de comparación, por lo que sus resultados pueden estar condicionados por el estilo y el contenido de las referencias disponibles: SkinVL, ajustado con las descripciones de \emph{MM-Skin}, obtiene en su propio conjunto de prueba y en una comparación no directa un BLEU-4 más de cien veces
mayor que el de los generadores de este trabajo con un BERTScore parecido (Sección~\ref{sec:res_gen_estado_arte}). La
consistencia diagnóstica, a su vez, favorece enumerar, y por eso LLaVA-Med obtiene sin anclaje la consistencia más alta
nombrando dos clases por reporte (Sección~\ref{sec:res_gen_sin_anclaje}). Conviene, entonces, que una métrica de fidelidad se acompañe de una medida de cuánto se compromete el reporte, como la proporción de reportes que nombran solo la clase verdadera.

Las comparaciones con el estado del arte muestran que el mayor uso de recursos computacionales no garantiza, por sí solo, un mejor desempeño bajo cualquier protocolo. Sin embargo, las diferencias entre trabajos impiden determinar cuánto de la distancia observada puede atribuirse específicamente a la escala. El \emph{pipeline} se ajusta en una sola GPU de 24~GB y su generación entra en una de consumo de 12~GB: Mini-VLM ocupa como máximo 8.6~GB en la A30 y 8.3~GB en la RTX 4070, y en la A30 escribe un reporte en unos 16 segundos, y el clasificador y el \gls{knn} agregan un costo mínimo (Secciones~\ref{sec:res_costo_clasificadores}
y~\ref{sec:res_costo}). Los trabajos revisados, en cambio, ajustan sus generadores en ocho GPU, preentrenan sus modelos
fundacionales en cuatro H100 o redactan con un modelo de 70 mil millones de parámetros consultado por
API~\cite{khan2025ensemble} (Tabla~\ref{tab:marco_generadores} del Capítulo~\ref{cp:marco}). En la clasificación, esa
diferencia no parece decisiva, porque la distancia con los mejores resultados de la literatura es del orden del efecto
de la partición, y entre los modelos evaluados el más grande no mejora a Swin-Tiny. En la representación contrastiva, en cambio, sí parece serlo: DermFM-Zero, preentrenado con tres millones de imágenes en cuatro H100 y alineado después con un millón de pares
imagen--texto, supera a DermLIP en el protocolo de
sus autores con 0.733 frente a 0.561~\cite{yan2026dermfmzero}. Entre los generadores evaluados, por último, el tamaño no ordena la
fidelidad: LLaVA-Med, con casi el doble de parámetros que Mini-VLM, transcribe peor el diagnóstico, y SkinVL, del mismo
tamaño que LLaVA-Med, describe la mitad de las entidades clínicas (Sección~\ref{sec:res_gen_comparacion}).

En las condiciones experimentales evaluadas, la restricción de utilizar una sola GPU con hasta 24~GB de memoria resultó compatible con la ejecución del \emph{pipeline} y con la consistencia diagnóstica obtenida. Este resultado se encuentra asociado a la configuración particular del clasificador y del generador seleccionados. Lo que la restricción deja fuera son dos
cosas que este trabajo no evaluó: los generadores de mayor tamaño, cuya descripción podría ser distinta, y el
preentrenamiento contrastivo a gran escala, que es la vía por la que la literatura mejora la representación que aquí
usa el \gls{knn}.

En conjunto, los resultados sugieren que el desempeño observado está condicionado por una combinación de factores. La disponibilidad y distribución de los datos pueden limitar la capacidad de aprendizaje de los clasificadores, las decisiones de evaluación influyen sobre las métricas obtenidas y las restricciones computacionales determinan qué modelos y procedimientos pueden estudiarse. Aunque los experimentos no permiten aislar la contribución de cada factor, muestran la importancia de considerar estas condiciones al interpretar los resultados y compararlos con los de la literatura.

\section{Alcance de los resultados y limitaciones}
\label{sec:disc_limitaciones}
% Extensión estimada: 3 páginas, dos listas y cinco párrafos (ficha 6.5).
% Fuentes: sec:met_eval_generacion, sec:met_estadistica, sec:met_referencias, sec:met_pipeline (piloto de la plantilla),
% sec:met_hiba; sec:res_gen_nivel1, sec:res_gen_ablaciones, sec:res_gen_cualitativos, sec:res_hiba, sec:res_robustez,
% sec:res_sup_clases, sec:res_gen_clases, sec:res_gen_pipeline, sec:res_costo; sec:estado_arte_reportes; secciones
% "Amenazas a la validez" y "Limitaciones" de Docs/experimentacion/experimentacion.tex, con las cifras de la partición
% por lesión. Recibe del Capítulo 1 las limitaciones de alcance. Reemplaza las antiguas 6.4 (métricas) y 6.7.

Las conclusiones de las secciones anteriores valen dentro de las condiciones en que se obtuvieron: imágenes
dermatoscópicas de \emph{HAM10000}, métricas automáticas y el cómputo de una sola GPU. La primera de esas condiciones
que conviene precisar es qué establecen las métricas. La consistencia diagnóstica busca el nombre de la clase en el
texto, de modo que mide la fidelidad diagnóstica del reporte y no si un dermatólogo lo aceptaría, y las métricas con
referencia del Nivel~1 miden sobre todo el parecido de estilo con un texto ajeno (Sección~\ref{sec:res_gen_nivel1}).
En cuanto a la descripción, se evaluó que cambia con la imagen (Sección~\ref{sec:res_gen_ablaciones}) y cuántas
entidades clínicas nombra. Comprobar que cada rasgo descrito esté en la imagen requiere la validación adicional de un
especialista, y los ejemplos de la Sección~\ref{sec:res_gen_pipeline} muestran defectos que estas métricas automáticas
no registran. Esta limitación es compartida por la literatura revisada, en la que la evaluación con
expertos es escasa: cinco dermatólogos en SkinGPT-4~\cite{zhou2024pretrained}, dos entre veinte participantes en
SkinSavvy2~\cite{kim2025skinsavvy2} y un médico en el trabajo de \textcite{chakrabarti2025skin}, mientras que los
demás trabajos evalúan el texto por su parecido con una referencia o por la exactitud del diagnóstico
(Sección~\ref{sec:estado_arte_reportes}).

Además de esa distancia entre lo medido y la validez clínica, conviene tener en cuenta cuatro aspectos de la validez de los resultados:
\begin{itemize}
    \item \textbf{Validez interna.} No todos los experimentos usan el mismo protocolo de inferencia. Los resultados
    principales de clasificación promedian dos vistas de cada imagen con \gls{tta}, mientras que la robustez y la
    interpretabilidad usan una sola vista, y el costo se midió en tres GPU distintas. Cada tabla declara su protocolo,
    pero los valores absolutos de tablas distintas no siempre son comparables. Además, las variantes de la cascada sin
    imagen y con un diagnóstico equivocado usan un subconjunto de 400 imágenes con Mini-VLM y LLaVA-Med
    (Sección~\ref{sec:met_referencias}), de modo que sus intervalos son más amplios que los del conjunto de prueba completo.
    \item \textbf{Validez de la conclusión estadística.} Las comparaciones estadísticas usan como unidad la imagen de prueba y se
    calculan sobre la ejecución de referencia de cada modelo (Sección~\ref{sec:met_estadistica}). Miden, por lo
    tanto, la variación debida a la muestra de imágenes, pero no la variación entre entrenamientos, cuyo rango supera una décima en las clases minoritarias y en el melanoma, y que se informa aparte como desvío entre las diez ejecuciones. La ejecución de referencia, fijada de antemano, queda por encima de la media de las diez, de modo que la cascada se mide con un clasificador algo mejor que el promedio: en melanoma, su sensibilidad es 0.701, frente a 0.585 de media. Las
    clases con pocas imágenes, como \gls{df} con 17, dan intervalos amplios, y algunas diferencias, como la del
    ensamble frente a Swin-Tiny o la de DermLIP con una cabeza frente a los supervisados, no se compararon estadísticamente.
    \item \textbf{Validez de constructo.} La consistencia diagnóstica no juzga el reporte, el Nivel~1 mide sobre todo
    el estilo, y la cobertura de entidades cuenta términos clínicos sin comprobar que correspondan a la imagen. En el
    nivel diferencial, además, la consistencia se parece a acertar entre las dos primeras opciones y no equivale a una
    exactitud.
    \item \textbf{Validez externa.} La evaluación principal usa un solo conjunto, y en \emph{HIBA} Swin-Tiny pierde
    cerca del 40\,\% de su sensibilidad macro, también en las imágenes dermatoscópicas, con un mapeo de diagnósticos
    que agrega ruido en \gls{akiec} (Sección~\ref{sec:res_hiba}). Como la fidelidad diagnóstica del reporte coincide
    con la exactitud del clasificador, es esperable que también caiga en otro conjunto de datos, aunque no se midió. Las
    comparaciones con la literatura, por su parte, ubican el orden de magnitud de los resultados, pero no permiten
    ordenar los modelos.
\end{itemize}

A estos aspectos se suman las limitaciones de alcance del trabajo. La principal es la ausencia de especialistas.
Ningún dermatólogo leyó los reportes, y no existe para las imágenes de \emph{HAM10000} un conjunto de reportes
redactados por especialistas con el que compararlos. Por eso el trabajo mide la fidelidad diagnóstica de los reportes, pero no su validez clínica. Tampoco se pronuncia sobre la fotografía clínica ni sobre la histopatología, donde lo medido no puede trasladarse sin volver a medirlo.

Por último, no todos los resultados tienen el mismo respaldo. Las decisiones de selección se tomaron en validación: el modelo, por la media de las ejecuciones, y el umbral de melanoma, con una especificidad mínima fijada de antemano. El conjunto de prueba de la partición por lesión y \emph{HIBA} aportan, entonces, una evaluación independiente de esas decisiones, y el ruteo no ajusta ningún umbral sobre el conjunto de prueba. La única excepción es la instrucción \emph{noreasoning}, que se eligió en un piloto sobre 100 imágenes, de las cuales 17 pertenecen al conjunto de prueba (Sección~\ref{sec:met_pipeline}), sin fijar ningún umbral sobre él. En cambio, las explicaciones que este capítulo propone, como la jerarquía entre las fuentes del reporte o el origen de los errores compartidos, y los análisis que se hicieron después de ver los
resultados, como la clasificación de las fallas de los reportes o los ejemplos, describen el conjunto de prueba sin haberse
confirmado en otros datos, y quedan como exploratorios. Con estas limitaciones delimitadas, el
Capítulo~\ref{cp:conclusiones} responde las preguntas de investigación y reúne las conclusiones del trabajo.

